\section{Geometry and Order Dependence of the Exact Threshold}\label{sec:geometry}

The exact criterion in Theorem~\ref{thm:c10} reduces the complete positive-diagonal orbit to the scalar threshold \(T_\alpha(\beta)\). We next show that this variational quantity has substantially more structure than a generic two-variable minimum.

\subsection{Strict Convexity and a Unique Optimizer}

Fix \(2/3<\alpha<1\). Let
\[
 u=\cos\!\left(\frac{\alpha\pi}{2}\right),\qquad
 K_\alpha=1-4u^2,
\]
and recall from Equation~\eqref{eq:halpha} that \(h_\alpha(b)=r^2(1+2ur)\), where \(r\) is the positive solution of
\[
 b=K_\alpha r^2-2ur.
\]
Introduce the elasticity
\begin{equation}
 E_\alpha(b)
 :=
 \frac{b\,h_\alpha'(b)}{h_\alpha(b)}.
 \label{eq:elasticity}
\end{equation}

\begin{Lemma}[elasticity bounds]\label{lem:elasticity}
For every \(b>0\),
\begin{equation}
 0<E_\alpha(b)<\frac32,
 \qquad
 \frac{dE_\alpha}{db}>0.
 \label{eq:elasticity-bounds}
\end{equation}
\end{Lemma}

\begin{proof}
Using \(r\) as parameter,
\[
 \frac{db}{dr}=2(K_\alpha r-u),
 \qquad
 \frac{dh_\alpha}{dr}=2r(1+3ur),
\]
and therefore
\begin{equation}
 E_\alpha
 =
 \frac{(K_\alpha r-2u)(1+3ur)}
      {(K_\alpha r-u)(1+2ur)}.
 \label{eq:E-r}
\end{equation}
On \(r>2u/K_\alpha\), both factors
\[
 \frac{K_\alpha r-2u}{K_\alpha r-u},
 \qquad
 \frac{1+3ur}{1+2ur}
\]
increase strictly with \(r\). Their product increases from \(0\) to \(3/2\). Since \(b(r)\) is strictly increasing, Equation~\eqref{eq:elasticity-bounds} follows.
\end{proof}

Parameterize the open simplex by logit variables
\begin{equation}
 x_1=\frac{e^{y_1}}{e^{y_1}+e^{y_2}+1},
 \quad
 x_2=\frac{e^{y_2}}{e^{y_1}+e^{y_2}+1},
 \quad
 x_3=\frac{1}{e^{y_1}+e^{y_2}+1}.
 \label{eq:logit}
\end{equation}
Set
\[
 L(y)=\log(e^{y_1}+e^{y_2}+1),
\]
\[
 Q_\beta(y)
 =
 \log\!\left(
 \beta_{12}e^{y_1+y_2}
 +\beta_{13}e^{y_1}
 +\beta_{23}e^{y_2}
 \right),
\]
and \(\phi_\alpha(s)=\log h_\alpha(e^s)\). Then
\begin{equation}
 G_{\alpha,\beta}(y)
 :=
 \log\!\left[
 \frac{h_\alpha(B_\beta(x(y)))}{x_1x_2x_3}
 \right]
 =
 \phi_\alpha(Q_\beta-2L)+3L-y_1-y_2.
 \label{eq:logobjective}
\end{equation}

\begin{Theorem}[strict threshold geometry]\label{thm:strict-convexity}
For every \(\beta\in\mathbb R_{>0}^3\) and \(2/3<\alpha<1\), \(G_{\alpha,\beta}\) is globally strictly convex on \(\mathbb R^2\). Hence:
\begin{enumerate}
\item \(T_\alpha(\beta)\) has a unique minimizer \(x^*(\alpha,\beta)\in\Delta_2^\circ\);
\item the minimizer is nondegenerate;
\item \(x^*\) and \(T_\alpha\) depend smoothly on \((\alpha,\beta)\).
\end{enumerate}
\end{Theorem}

\begin{proof}
Put \(s=Q_\beta-2L\). Since
\[
 \phi_\alpha'(\log b)=E_\alpha(b),
 \qquad
 \phi_\alpha''(\log b)=b\,E_\alpha'(b),
\]
Lemma~\ref{lem:elasticity} gives
\[
 \phi_\alpha'>0,\qquad
 \phi_\alpha''>0,\qquad
 3-2\phi_\alpha'>0.
\]
Differentiating Equation~\eqref{eq:logobjective},
\begin{equation}
 \nabla^2G_{\alpha,\beta}
 =
 \phi_\alpha''\,\nabla s\nabla s^\top
 +\phi_\alpha'\nabla^2Q_\beta
 +(3-2\phi_\alpha')\nabla^2L.
 \label{eq:hessian-decomp}
\end{equation}
Both \(Q_\beta\) and \(L\) are log-sum-exp functions of affine forms, so \(\nabla^2Q_\beta\succeq0\); moreover \(\nabla^2L\succ0\) because all three softmax weights are positive. Every coefficient in Equation~\eqref{eq:hessian-decomp} is nonnegative and the coefficient of \(\nabla^2L\) is strictly positive. Thus \(\nabla^2G_{\alpha,\beta}\succ0\) everywhere.

As \(\|y\|\to\infty\), \(x(y)\) approaches the simplex boundary and \(x_1x_2x_3\to0\), while \(h_\alpha(B_\beta(x))\) remains positive. Hence \(G_{\alpha,\beta}\) is coercive. A unique minimizer exists; positive definiteness gives nondegeneracy, and the implicit-function theorem gives smooth dependence.
\end{proof}

Two exact reductions follow immediately. If \(\beta_{12}=\beta_{13}=\beta_{23}=b\), uniqueness and permutation symmetry imply \(x^*=(1/3,1/3,1/3)\), so
\begin{equation}
 T_\alpha(b,b,b)=27h_\alpha(b/3).
 \label{eq:symmetric-slice}
\end{equation}
For \(b=1\), this becomes the cyclic-network boundary
\begin{equation}
 T_\alpha(1,1,1)=1+R_3(\alpha)^3,
 \qquad
 R_3(\alpha)
 =
 \frac{\sin(\alpha\pi/2)}
 {\sin(\alpha\pi/2-\pi/3)},
 \label{eq:siami-slice}
\end{equation}
recovering the structured slice studied by Siami \cite{Siami2021}. If two \(\beta\)'s are equal, the corresponding two simplex coordinates of \(x^*\) are equal, reducing the minimization exactly to one dimension.

\subsection{Dependence on the Fractional Order}

\begin{Theorem}[monotone widening and the classical limit]\label{thm:alpha-deformation}
For fixed \(\beta>0\), \(T_\alpha(\beta)\) is strictly decreasing in \(\alpha\) on \(2/3<\alpha\le1\). Let
\[
 S=\sqrt{\beta_{12}}+\sqrt{\beta_{13}}+\sqrt{\beta_{23}},
 \qquad
 G=\sqrt{\beta_{12}\beta_{13}\beta_{23}}.
\]
As \(\alpha\to1^{-}\),
\begin{equation}
 T_\alpha(\beta)
 =
 T_1(\beta)
 +
 C(\beta)(1-\alpha)
 +
 O((1-\alpha)^2),
 \label{eq:classical-rate}
\end{equation}
where
\begin{equation}
 C(\beta)
 =
 \pi\left(
 \frac{S^{5/2}}{\sqrt G}
 +
 S^{3/2}\sqrt G
 \right)>0.
 \label{eq:Cbeta}
\end{equation}
\end{Theorem}

\begin{proof}
For fixed positive \(a,b\), decreasing the boundary angle \(\theta=\alpha\pi/2\) strictly enlarges the Matignon region. The unique coefficient value \(H_\theta(a,b)\) placing a conjugate pair on the boundary therefore increases strictly as \(\theta\) decreases. Hence \(h_{\alpha_1}(b)>h_{\alpha_2}(b)\) whenever \(\alpha_1<\alpha_2\), and minimization over the same simplex preserves strictness.

For the rate, \(u=\cos(\alpha\pi/2)=\frac{\pi}{2}(1-\alpha)+O((1-\alpha)^3)\). Implicit differentiation of
\[
 (1-4u^2)r^2-2ur-b=0
\]
at \(u=0\) gives \(\partial r/\partial u=1\). Consequently
\begin{equation}
 h_\alpha(b)
 =
 b+\pi(1-\alpha)\sqrt b(1+b)+O((1-\alpha)^2).
 \label{eq:h-expansion}
\end{equation}
At \(\alpha=1\), the unique simplex minimizer is
\[
 x_1^*=\frac{\sqrt{\beta_{23}}}{S},
 \quad
 x_2^*=\frac{\sqrt{\beta_{13}}}{S},
 \quad
 x_3^*=\frac{\sqrt{\beta_{12}}}{S},
\]
for which
\[
 x_1^*x_2^*x_3^*=\frac{G}{S^3},
 \qquad
 B_\beta(x^*)=\frac{G}{S}.
\]
The envelope theorem applied to the nondegenerate minimum and Equation~\eqref{eq:h-expansion} yields Equations~\eqref{eq:classical-rate}--\eqref{eq:Cbeta}.
\end{proof}

Theorem~\ref{thm:alpha-deformation} quantifies the part of the order dependence that is needed below: the fractional-only band widens strictly as \(\alpha\) decreases and closes linearly at the classical endpoint. Together with the exact symmetric reductions above, this supplies a smooth and uniquely optimized threshold geometry without adding further parameters to the C-10 classification.
